% paper/sections/10_1_spatial.tex
%
% §10.1  空間離散化の精度検証
%   CCD/DCCD 収束性・散逸特性・曲率精度・ψ直接法・ε_eff 診断
% ═══════════════════════════════════════════════════════════════════════════════
\subsection{空間離散化の精度検証}
\label{sec:verify_spatial}
% ═══════════════════════════════════════════════════════════════════════════════

本稿の全離散演算子は CCD 法を基盤としている．
本節ではまず CCD 微分演算子そのものの精度（§\ref{sec:verify_ccd_convergence}）を確認し，
次にその拡張である DCCD フィルタの散逸特性（§\ref{sec:verify_dccd_dissipation}，§\ref{sec:dccd_dissipation_bench}），
曲率計算への応用精度（§\ref{sec:verify_curvature}），
非一様格子への適用性（§\ref{sec:gcl_verification}）を順に検証する．

% ─────────────────────────────────────────────────────────────────────────────
\subsubsection{CCD/DCCD 微分演算の収束性}
\label{sec:verify_ccd_convergence}
% ─────────────────────────────────────────────────────────────────────────────

\noindent\textbf{目的：}
CCD 法（§\ref{sec:ccd_def}）が設計精度 $\Ord{h^6}$ を達成することを格子収束テストで確認する．

\noindent\textbf{評価対象：}
CCD 1 階微分および 2 階微分演算子（周期境界条件）．

\noindent\textbf{手順とパラメタ：}
テスト関数 $f = \sin(2\pi x)\sin(2\pi y)$（解析微分既知）を用い，
$N = 16, 32, 64, 128, 256$ で格子を精緻化し，$L_2$・$L_\infty$ 誤差と収束次数 $p$ を計測する．
境界条件は周期．

\noindent\textbf{結果：}

\begin{table}[htbp]
\centering
\caption{CCD 1 階微分の格子収束（周期境界，$f = \sin(2\pi x)\sin(2\pi y)$）}
\label{tab:ccd_periodic}
\renewcommand{\arraystretch}{1.3}
\begin{tabular}{rrrrr}
\toprule
$N$ & $L_2$ 誤差 & 次数 & $L_\infty$ 誤差 & 次数 \\
\midrule
 16 & $1.28 \times 10^{-6}$  & ---  & $2.57 \times 10^{-6}$  & ---  \\
 32 & $1.93 \times 10^{-8}$  & 6.06 & $3.86 \times 10^{-8}$  & 6.06 \\
 64 & $2.99 \times 10^{-10}$ & 6.01 & $5.97 \times 10^{-10}$ & 6.01 \\
128 & $4.65 \times 10^{-12}$ & 6.00 & $9.36 \times 10^{-12}$ & 6.00 \\
256 & $7.67 \times 10^{-14}$ & 5.92 & $2.28 \times 10^{-13}$ & 5.36 \\
\bottomrule
\end{tabular}
\end{table}

1 階微分は $N = 16 \to 128$ で $p \approx 6.00$--$6.06$ の収束を示した．
2 階微分も同等の $\Ord{h^{5.97}}$ 精度を達成した．
$N = 256$ では打ち切り誤差が $\Ord{10^{-14}}$ に到達し，
倍精度浮動小数点の丸め誤差により見かけ上の収束次数が低下した
（$L_\infty$ で $p \approx 5.36$）．

\noindent\textbf{結論：}
CCD 微分演算子は設計精度 $\Ord{h^6}$ を達成した．

% ─────────────────────────────────────────────────────────────────────────────
\subsubsection{DCCD の数値散逸特性}
\label{sec:verify_dccd_dissipation}
% ─────────────────────────────────────────────────────────────────────────────

\noindent\textbf{目的：}
DCCD フィルタ（§\ref{sec:dccd_decoupling}）が Nyquist モード（チェッカーボード振動）を
設計通りに抑制することを解析的・数値的に確認する．

\noindent\textbf{評価対象：}
DCCD 離散伝達関数 $H(\xi; \varepsilon_d)$．

\noindent\textbf{手順とパラメタ：}
伝達関数を $e^{i\xi x}$ に対する CCD の応答比として定義し，
解析式
\begin{equation}
  H(\xi; \varepsilon_d) = 1 - 4\varepsilon_d \sin^2(\xi/2)
  \label{eq:dccd_transfer_verify}
\end{equation}
を $\varepsilon_d = 0.00, 0.05, 0.25$ で評価する．
$N = 128$ 格子上でチェッカーボードモード $(-1)^{i+j}$ に
$\varepsilon_d = 0.25$ を適用し，フィルタ後の振幅を数値的に確認する．

\noindent\textbf{結果：}

\begin{table}[htbp]
\centering
\caption{DCCD 伝達関数の Nyquist 波数における値 $H(\pi; \varepsilon_d)$}
\label{tab:dccd_transfer}
\renewcommand{\arraystretch}{1.3}
\begin{tabular}{rr}
\toprule
$\varepsilon_d$ & $H(\pi)$ \\
\midrule
0.00 & $1.0000$（散逸なし：CCD）\\
0.05 & $0.8000$（移流用：低散逸）\\
0.25 & $0.0000$（PPE 用：完全抑制）\\
\bottomrule
\end{tabular}
\end{table}

$\varepsilon_d = 1/4$ で $H(\pi) = 0$（§\ref{sec:checkerboard_solution} の理論と一致）．
$\varepsilon_d = 0.05$（移流用）では $H(\pi) = 0.80$ であり，
高波数成分を 20\% 減衰させつつ低波数の移流精度を保持する．
$N = 128$ での数値実験でもフィルタ後のチェッカーボード振幅がゼロとなることを確認した．

\noindent\textbf{結論：}
DCCD フィルタは設計通り Nyquist モードを完全に抑制し，低波数成分を保存する．

% ─────────────────────────────────────────────────────────────────────────────
\subsubsection{\texorpdfstring{曲率 $\kappa$ の計算精度}{曲率 κ の計算精度}}
\label{sec:verify_curvature}
% ─────────────────────────────────────────────────────────────────────────────

\noindent\textbf{目的：}
曲率計算（§\ref{sec:curvature}，式~\eqref{eq:curvature_2d}）の 3 つの計算経路を比較し，
本稿で採用する経路 B の精度特性を確認する．

\noindent\textbf{評価対象：}
3 つの曲率計算経路：
\begin{itemize}[noitemsep]
  \item \textbf{経路 A（$\phi$ 直接）：}
        真の符号付距離関数 $\phi$ に CCD を適用し $\kappa$ を計算（理想ベースライン）．
  \item \textbf{経路 B（$\psi \to \phi$ ロジット逆変換）：}
        CLS 変数 $\psi$ からロジット逆変換
        $\phi = \varepsilon\ln[\psi/(1-\psi)]$（§\ref{sec:newton_inversion}）で
        $\phi$ を復元し，$\phi$ の CCD 微分から $\kappa$ を計算（採用経路）．
  \item \textbf{経路 C（$\psi$ 直接法，§\ref{sec:curvature_direct_psi_cls}）：}
        $\psi$ の CCD 微分値から式~\eqref{eq:curvature_psi_2d} で $\kappa$ を直接計算．
\end{itemize}

\noindent\textbf{手順とパラメタ：}
2 種類のテスト問題を用いる．
\begin{itemize}[noitemsep]
  \item[(a)] 正弦波状界面 $y = 0.5 + 0.05\sin(2\pi x)$：
        解析解 $\kappa = f''/(1+f'^2)^{3/2}$，$\varepsilon = 1.5h$，$N = 32$--$256$．
  \item[(b)] 円形界面 $R = 0.25$：$\kappa_\mathrm{exact} = 4$．
\end{itemize}
評価指標は $L_\infty$ 誤差と収束次数．

\noindent\textbf{結果：}

\begin{table}[htbp]
\centering
\caption{正弦波状界面の曲率 $L_\infty$ 誤差収束（$A = 0.05$，$\varepsilon = 1.5h$）}
\label{tab:curvature_sinusoidal}
\renewcommand{\arraystretch}{1.3}
\begin{tabular}{rrrrrrr}
\toprule
$N$ & A: $\phi$ 直接 & 次数 & B: ロジット経由 & 次数 & C: $\psi$ 直接 & 次数 \\
\midrule
 32 & $7.64 \times 10^{-4}$ & ---  & $7.64 \times 10^{-4}$ & ---  & $3.84 \times 10^{-2}$ & ---  \\
 64 & $2.45 \times 10^{-5}$ & 4.96 & $2.14 \times 10^{-5}$ & 5.16 & $2.19 \times 10^{-2}$ & 0.81 \\
128 & $7.72 \times 10^{-7}$ & 4.99 & $5.80 \times 10^{-6}$ & 1.88 & $2.27 \times 10^{-2}$ & $-0.05$ \\
256 & $2.42 \times 10^{-8}$ & 5.00 & $1.33 \times 10^{-5}$ & $-1.20$ & $3.59 \times 10^{-2}$ & $-0.66$ \\
\bottomrule
\end{tabular}
\end{table}

経路 A は全格子で $\Ord{h^5}$ を維持した．
経路 B は $N \leq 64$ で A と同等の精度を示したが，$N \geq 128$ では
ロジット逆変換の飽和域処理に起因する誤差が顕在化し収束が停滞した．
経路 C は全格子で収束しなかった．
円形界面（$R = 0.25$）では経路 B の収束は $\Ord{h^1}$ にとどまり，
界面厚み $\varepsilon = \Ord{h}$ が精度の上界を規定した．

\noindent\textbf{結論：}
本稿では経路 B を採用する．
GFM 界面処理（§\ref{sec:gfm}）の精度が $\Ord{h^2}$（界面近傍）で支配的であるため，
曲率計算経路の選択は全体精度に影響しない．
なお，曲率の計算式（式~\eqref{eq:curvature_2d}）は
$\kappa = (\phi_{xx}\phi_y^2 - 2\phi_{xy}\phi_x\phi_y + \phi_{yy}\phi_x^2)/|\bnabla\phi|^3$
であり，混合微分 $D_{xy}^{(11)}$ を含む．
上記の曲率収束テストで $\Ord{h^5}$ が確認されたことは，
$D_{xy}^{(11)}$（2 方向 CCD の逐次適用）が設計精度を発揮していることの
間接的な検証となっている．

% ─────────────────────────────────────────────────────────────────────────────
\subsubsection{非一様格子上の微分精度と幾何学的保存則（GCL）}
\label{sec:gcl_verification}
% ─────────────────────────────────────────────────────────────────────────────
\input{sections/10_1a_gcl_verification}

% ─────────────────────────────────────────────────────────────────────────────
\subsubsection{DCCD フィルタ：1次元移流問題による散逸特性の定量評価}
\label{sec:dccd_dissipation_bench}
% ─────────────────────────────────────────────────────────────────────────────
% paper/sections/10_1b_dccd_dissipation.tex
%
% §10.1.5  DCCD フィルタ：分散・散逸特性と選択的高周波減衰の評価
%          実験スクリプト: experiments/dccd_comparison.py
%          結果データ:    results/dccd_comparison/

\noindent\textbf{目的：}
DCCD の 10 次選択フィルタ（§~\ref{sec:advection_dccd_design}）が，
不連続プロファイルに対して CCD の過大な全変動（TV）を大幅に低減しつつ，
滑らかなプロファイルでは CCD と同等の高精度を維持することを検証する．

\noindent\textbf{評価対象：}
DCCD の分散・散逸特性．CLS 法の界面 $\psi$（準不連続プロファイル）に対する
DCCD の選択的高波数減衰の有効性を間接的に評価する．

\noindent\textbf{手順とパラメタ：}
1次元線形移流方程式
\begin{equation}
  \frac{\partial u}{\partial t} + c\,\frac{\partial u}{\partial x} = 0,
  \qquad x \in [0,1),\quad c = 1,\quad T = 1
  \label{eq:linear_advection_dccd}
\end{equation}
を以下の条件で解く：
\begin{itemize}[noitemsep]
  \item 格子 $N=256$，周期境界，CFL $=0.4$，時間積分 RK4，1 周期（$T=1$）
  \item 初期条件 3 種：Square（不連続矩形パルス），Triangle（$C^0$ 三角波），Smooth tanh
  \item 比較スキーム：O2（2次中心差分），O4（4次），CCD（6次），
        DCCD（CCD $+$ 10次選択フィルタ，$\alpha_f=0.4$），WENO5
  \item 評価指標：$L_2 = \sqrt{\mathrm{mean}(u-u_\mathrm{exact})^2}$，
        $\mathrm{TV} = \sum_i |u_{i+1}-u_i|$
\end{itemize}

\noindent\textbf{結果：}

\begin{table}[htbp]
\centering
\caption{1周期後の $L_2$ 誤差（$N=256$，CFL$=0.4$，太字：各初期条件での最良値）}
\label{tab:dccd_l2_bench}
\begin{tabular}{lrrrrr}
\toprule
初期条件 & O2 & O4 & CCD & DCCD & WENO5 \\
\midrule
Square（不連続） & $1.41\times10^{-1}$ & $9.08\times10^{-2}$ & $4.80\times10^{-2}$ & $\mathbf{4.23\times10^{-2}}$ & $6.35\times10^{-2}$ \\
Triangle（$C^0$） & $2.41\times10^{-2}$ & $5.84\times10^{-3}$ & $\mathbf{1.37\times10^{-3}}$ & $1.41\times10^{-3}$ & $3.90\times10^{-3}$ \\
Smooth tanh & $3.75\times10^{-2}$ & $2.46\times10^{-3}$ & $\mathbf{2.57\times10^{-5}}$ & $3.75\times10^{-5}$ & $9.66\times10^{-4}$ \\
\bottomrule
\end{tabular}
\end{table}

\begin{table}[htbp]
\centering
\caption{1周期後の全変動（TV）．
  Exact 値は初期条件の TV（Square: 2.000, Triangle: 1.969, tanh: 2.000）．
  太字：各初期条件での最良値（TV $\to$ exact に最近）．}
\label{tab:dccd_tv_bench}
\begin{tabular}{lrrrrr}
\toprule
初期条件 & O2 & O4 & CCD & DCCD & WENO5 \\
\midrule
Square（不連続） & $18.40$ & $17.17$ & $10.83$ & $3.15$ & $\mathbf{2.00}$ \\
Triangle（$C^0$） & $2.60$ & $2.37$ & $2.09$ & $1.97$ & $\mathbf{1.91}$ \\
Smooth tanh & $2.82$ & $2.02$ & $2.00$ & $2.00$ & $\mathbf{2.00}$ \\
\bottomrule
\end{tabular}
\end{table}

\noindent\textbf{結論：}
\begin{itemize}[noitemsep]
  \item \textbf{不連続（Square）：} DCCD は CCD に比べ TV を $10.83 \to 3.15$ に低減（約 $1/3$）し，
        $L_2$ 誤差も最良（$4.23\times10^{-2}$）であった．
  \item \textbf{平滑（Smooth tanh）：} DCCD と CCD の $L_2$ 誤差は同オーダー
        （$3.75\times10^{-5}$ vs $2.57\times10^{-5}$，1.5 倍以内），TV も同値（$2.000$）．
        フィルタが高波数成分のみを対象とし，滑らかな成分を保護することを確認した．
  \item \textbf{CLS 法への含意：} CLS の界面 $\psi$ は Smooth tanh に相当するため，
        DCCD は CCD と同等の精度で界面を移流しつつ振動を抑制できる．
        §~\ref{sec:interface_advection_bench}（Zalesak・単一渦）の低体積誤差の理論的根拠となる．
\end{itemize}

